\section{问题二的模型建立与求解}

\subsection{多点多架次调度模型}

问题二在问题一的航段几何与能耗模型基础上，放开单点往返限制：每个架次可依次访问一个或多个服务区后返回 $O01$，并引入实体机队（A/B/C 三型共 8 架）、共享电池数量、两阶段充电周转及物资时限约束。设架次集合为 $\mathcal{R}$，每个架次 $r$ 需确定其货箱集合、访问顺序 $\pi_r=(O01,S_{r_1},\dots,S_{r_{n_r}},O01)$、机型 $g_r$、执行无人机 $u_r$、电池 $c_r$ 及开始时刻 $s_r$。

\textbf{载荷与容量约束.} 架次在每一航段的剩余载荷须满足对应机型的能量/质量与体积约束。设架次沿途投送后剩余载荷序列为 $\{q_r^{(\ell)}\}$，则各航段能耗按问题一公式累加，且须满足返航安全余量：
\begin{equation}
\sum_{\ell} E^{(\ell)}_r\big(q_r^{(\ell)}\big)\;\le\;(1-\eta_{g_r})\,B_{g_r}.
\end{equation}

\textbf{资源与充电约束.} 同一机型的共享电池可在该机型不同实体机间调度，不同机型不可混用。电池任务结束剩余 SOC 按可用能量与实际能耗折算，不得低于返航电量下限；再次投入前须充至 $100\%$，充电时间采用两阶段（快充/涓流）等效模型\cite{ref_ti}：
\begin{equation}
T_{\text{chg}}(\text{SOC})=
\begin{cases}
\dfrac{0.9-\text{SOC}}{0.9}\cdot 0.65\,T^{\text{full}}+0.35\,T^{\text{full}}, & \text{SOC}<0.9,\\[2mm]
\dfrac{1-\text{SOC}}{0.1}\cdot 0.35\,T^{\text{full}}, & \text{SOC}\ge 0.9,
\end{cases}
\end{equation}
同一资源的任务占用与充电时段不得重叠，不同资源可并行充电。

\textbf{时限约束.} 医疗物资须在期望送达时间前送达，首批保障货箱须在首批截止时间前送达，其余物资的期望送达时间用于度量及时性（软约束）。

\textbf{优化目标.} 采用四目标：加权迟到量 $f_1$（以应急优先系数加权的迟到时间之和）、完工时间 $f_2$（所有无人机完成最后架次返回 $O01$ 的最晚时刻）、总能耗 $f_3$、架次数 $f_4$：
\begin{equation}
\min\ \big(f_1,\,f_2,\,f_3,\,f_4\big).
\end{equation}

\subsection{多方案求解与比较}

\textbf{方法一（全局分组编码 NSGA-II）.} 决策空间同时含组批、排序、机型与资源指派，高维强耦合，难以精确求解。采用分组编码遗传算法：染色体编码“货箱$\to$架次”的分组与架次内访问顺序，由结构保证 $80/80$ 全覆盖；解码时贪心指派机型与可用电池并推进离散事件仿真计算四目标；以 NSGA-II 的非支配排序与拥挤度选择维护 Pareto 前沿。该方案是问题二的全局高复杂度候选，正式种群 $80$、代数 $80$。

代表解按代码实际规则从非支配前沿中优先选取架次数最少、再取能耗较小者：$f_1=6.37\times10^{5}$、$f_2=16683.6\,$s（$4.63\,$h）、$f_3=82.777\,$kWh、$N=29$（其中多停靠架次 $7$ 个、占 $24.1\%$），硬约束违反数为 $0$；最终非支配前沿含 $11$ 个零违反解。

\textbf{方法二（可行路线集合/贪心合并基线）.} 以问题一单点集合划分方案（$N_0=18$）为起点，按同机型、时窗相容和容量可容原则合并邻近服务区架次，共合并 $2$ 轮，架次降至 $16$，总能耗降至 $58.179\,$kWh。该方法作为可解释的基础构造，不追求全局 Pareto 最优；其结果为 $f_1=2.01\times10^{5}$、$f_2=10173.7\,$s、$f_3=58.179\,$kWh、$N=16$，硬约束违反数为 $5$。

\textbf{方法三（路线构造＋可行性修复混合方案）.} 代码包 \texttt{p2-1} 先按相邻服务区构造并合并路线，再用统一资源仿真检验；初始路线若出现硬时限违反，则调用 \texttt{q2-1} 全局分组搜索修复路线集合。实际输出为 $N=30$、$f_1=7.36\times10^{5}$、$f_2=17328.9\,$s、$f_3=78.118\,$kWh、违反数 $0$。由于其可行性修复复用了 NSGA-II 路线搜索，此结果应视为混合流程，而非独立的纯贪心解。

三种方案比较见表~\ref{tab:q2_methods}。全局 NSGA-II 与混合修复方案均满足硬时限，前者架次少 $1$ 且完工更快，后者能耗低 $4.659\,$kWh；贪心基线架次和能耗最低，但有 $5$ 个违反。上述指标对应不同目标权衡，不能只按单一指标断定一种方法全面占优。

\begin{table}[htbp]
\centering
\caption{问题二三种实现方案对比}
\label{tab:q2_methods}
\renewcommand{\arraystretch}{1.2}
\small
\begin{tabular}{c|c|c|c|c|c}
\toprule
方案 & 方法定位 & 架次 & 完工时间/s & 总能耗/kWh & 硬约束违反数 \\
\midrule
全局方案 & 分组编码 NSGA-II & 29 & 16683.6 & 82.777 & 0 \\
混合方案 & 路线修复混合 & 30 & 17328.9 & 78.118 & 0 \\
基础方案 & 贪心合并 & 16 & 10173.7 & 58.179 & 5 \\
\bottomrule
\end{tabular}
\end{table}

\begin{figure}[htbp]
  \centering
  \subfloat[NSGA-II 四目标 Pareto 前沿]{\includegraphics[width=0.48\textwidth]{q2_1_01.png}}
  \hfill
  \subfloat[代表解调度甘特图]{\includegraphics[width=0.48\textwidth]{q2_1_02.png}}
  \\
  \subfloat[资源占用时间线]{\includegraphics[width=0.48\textwidth]{q2_1_03.png}}
  \hfill
  \subfloat[NSGA-II 与贪心基线对照]{\includegraphics[width=0.48\textwidth]{q2_2_01.png}}
  \caption{问题二多点多架次调度结果}
  \label{fig:q2}
\end{figure}

\subsection{资源使用与可行性检验}

三种方案的输出均只调用 4 架无人机——B 型 U05/U06、C 型 U07/U08，A 型 U01--U04 闲置。原因是 A 型载质量（$25\,$kg）与航程包络在本需求结构下不占优，重载服务区由 C 型承担、其余由 B 型覆盖更省架次。C 型为瓶颈资源（仅 2 架机身、4 组电池）；在全局 NSGA-II 代表解中，C 型架次数为 $10$，电池 C-01\textasciitilde04 循环飞行/充电，未出现 SOC 越限或电池冲突。三种方案均完成 $80/80$ 个货箱配送。典型多停靠架次如 $T_{015}:O01\to S_{007}\to S_{003}\to S_{015}\to O01$（C 型、3 停靠）。

\subsection{调度结构的多维分析}

对 NSGA-II 代表解的架次结构进一步剖析如图~\ref{fig:adv_q2}。图~\ref{fig:adv_q2}(a) 将四目标 Pareto 前沿投影到“完工时间—总能耗”平面并以颜色编码架次数，可见前沿呈现明显的“时间—能耗”反向权衡：完工越快的解往往架次更多、能耗更高，代表解位于前沿膝部。图~\ref{fig:adv_q2}(b) 的小提琴图显示 C 型单架次时长的中位数与离散度均高于 B 型，反映其承担了更远、更重的多停靠任务。图~\ref{fig:adv_q2}(c) 的堆叠柱刻画停靠点数构成：B 型以单点架次为主，C 型则含较多 $2\!\sim\!3$ 停靠架次，是多点合并降本的主力。图~\ref{fig:adv_q2}(d) 进一步表明单架次时长随停靠点数增加而系统性拉长，量化了“合并降架次”与“单架次变长”之间的取舍。

\begin{figure}[htbp]
  \centering
  \includegraphics[width=0.96\textwidth]{adv_q2.png}
  \caption{多点多架次调度结构分析}
  \label{fig:adv_q2}
\end{figure}
